\documentclass{article}
	\usepackage[UTF8]{ctex}
	\usepackage{amsmath}
	\usepackage{graphicx}
	\usepackage{underscore}
	\usepackage{geometry}
	\usepackage{anysize}
	\date{}
	\title{全新比赛机制题解}
	\geometry{top=10mm,bottom=15mm,left=25mm,right=25mm}
	
	\begin{document}
		\maketitle \vspace{-50pt}
		
		\section{dfs}
		
			\subsection{优化}
			
				\subsubsection{use +/-}
				\indent one：$3^{13}*13 = 20726199$\\
				\indent +/-：$2*(3+3^{2}+...+3^{12})+3^{13} = 2391483$
				
				\subsubsection{剪枝}
				\indent 若当前使用时间超过当前完成最短时间，则结束。但判断的时间也不少。这个方法有利有弊。实际测试中，运行时间减少。
				
		\section{for循环}
		\indent 如果不会dfs的话，用for循环。\\
		\indent 考验毅力的时候了！ \\
		\indent 快速修改代码方法 $Ctrl+R$

		\section{dp}
		\subsection{$O(n^{2})$}
		\indent $f[x][y]$代表当前第一位同学花费x时间，第二位同学花费y时间，第三位同学花费的最小时间，操作次数为
		$
		\sum_{i=1}^{13} t_{ij}*
		\sum_{i=1}^{13} t_{ij}
		*13*3=260*260*13*3=2636400
		$。
		
		\subsection{$O(n^{2})$}
		\indent $f[x][y][z]$代表是否存在当前第一位同学花费x时间，第二位同学花费y时间，第三位同学花费z时间的方案，最大操作次数为
		$
		\sum_{i=0}^{12}(1+20*i)^{3}=148370079
		$。\\
		\indent 实际时间复杂度小于此，因为采用的是随机化数据。如果以平均值计算，操作次数为
		$
		\sum_{i=0}^{12}(1+10*i)^{3}=18844059
		$
		 ，正常来说不会超时。\\
		\indent 存储方面，采用滚动数组，即$f[x][y][z][2]$，上一个状态和当前状态为0和1，每次两者的值进行交换。本题没有卡空间，故没必要优化空间，避免运行时间的增加。
		
	\end
	{document}